\documentclass[11pt]{article}
\usepackage{fontspec}
\setmainfont{Times New Roman}
\setsansfont{Arial}
\setmonofont{Consolas}
\usepackage{amsmath,amssymb,mathtools}
\usepackage{geometry}
\usepackage{microtype}
\usepackage{enumitem}
\geometry{a4paper,margin=2.28cm}
\setlength{\parindent}{1.15em}
\setlength{\parskip}{0pt}
\makeatletter
\renewcommand\footnotesize{\@setfontsize\footnotesize{7.5}{8.5}\selectfont}
\makeatother
\emergencystretch=120em
\allowdisplaybreaks
\sloppy
\hbadness=10000
\vbadness=10000
\hfuzz=2pt
\vfuzz=10pt
\raggedbottom

\begin{document}

% Унит: Папер33 цомплете в001. Дирецт Интерславиц Латин транслатион фром рефресхед Зенодо R122 Папер32/Папер33 соурце слице, Папер33 линес 525--569 / сегментс P33-С0001--P33-С0010. The олдер лоцал Папер33 финал-аудитед слице ис суперседед фор тис унит бецаусе ит омитс the R122/соурце-сцан опенинг of the тво-проблем параграпх.
\section*{33. Хиперкомплексне величины и теорија представјень в аритметичном поиманьу}
\emph{Atti Congresso Bologna 2 (1928), s. 71--73.}

\begin{center}
\textsc{Emmy Noether} (Готтинген -- Германиа)\\[-0.1em]
\rule{3em}{0.4pt}\\[0.8em]
\textsc{Хиперкомплексне величины и теорија представјень}\\
\textsc{в аритметичном поиманьу}\footnote{Подробно изложенье буде јавити се в \emph{Math. Zeitschr.}; в продолженьу там буде разсмотрјена такоже некомутативна Галоисова теорија. Ср. такоже о “разпадных польах”: R. Brauer и E. Noether, \emph{Ber. Berl. Ak.} 1927, а о обчих структурных теоремах -- доклад Г. Köthe в тих \emph{Atti}.}
\end{center}

Хочу показати, како се теорија представјень -- особито представјеньа груп -- може разумети како теорија класов модулов и идеалов; и како се те последнье знову подчиньејут обобченому појму групы, именно групам с операторами.\footnote{Групы с операторами изходе од W. Krull (\emph{Math. Zeitschr.} 23) и О. Schmidt (\emph{Math. Zeitschr.} 29). И тут дано обработанье проблема редукције даного представјеньа изходи -- в ограниченьу на комутативне тела представјеньа -- од W. Krull (\emph{Theorie und Anwendung der verallgemeinerten Abelschen Gruppen}, Heidelberger Ber. 1926).}

Под \emph{представјеньем групы} -- в даном польу $P$ -- разумеје се, како ведомо, следујуче: елементам $a,b,\ldots$ групы $G$ придружајут се матрице $A,B,\ldots$ с коефициентами из $P$ тако, же продукту $ab$ одговарја продукт $AB$; систем матриц ставаје се хомоморфным образом групы. С тым представјеньем дано јест једновременно представјенье “групного колца”, кторо состоји из всих “групных чисел”, т. ј. всих линеарных комбинациј $a\alpha+b\beta+\cdots+c\gamma$, кде $a,b,\ldots,c$ пробегајут все комбинације из конечно многих групных елементов, а $\alpha,\beta,\ldots,\gamma$ соответно из конечно многих елементов из $P$; при том $a,b,\ldots,c$ сматрајут се линеарне независными, а множенье јест дано групным множеньем и рачунскыми законами. Представјенье групного колца јест једновременно хомоморфно взходно к сложеньу; тако представјеньа груп подчиньајут се обчему проблему \emph{представјеньа колц}, кде потребна јест хомоморфност взходно к сложеньу и множеньу.

Тут сут суштна два проблема. Првы јест проблем \emph{редукције даного} представјеньа и пытанье о једнозначности иредуцибилных составных чести, кторе при том наставајут. Други состоји в пытаньу о целости можливых представјень -- може быти само иредуцибилных -- даного колца в даном, не нужно комутативном телу.

Првы проблем јест много простши; то се види уже из того, же ни о колцу, кторо треба представити, ни о -- в обче некомутативном -- телу представјеньа не сут потребне никаке специалне предпоставјеньа. Факт, же при представјеньу иде о матрице фиксованого степена, даје все потребне предпоставјеньа конечности. Проблем можно пополно обработати посредством теорије груп с операторами, на ктору зато кратко прејду.

Група $\mathfrak{M}$ называје се \emph{група с областју операторов}, ако једновременно даны јест систем симболов $\Theta,H,\ldots$ -- операторов -- тако, же комбинације $\Theta(a),H(a),\ldots$ за всако $a$ из $\mathfrak{M}$ сут једнозначно определьене, дајут елемент из $\mathfrak{M}$ и задовольајут дистрибутивны закон: $\Theta(ab)=\Theta(a)\Theta(b)$. Две групы с једнакоју областју операторов называјут се операторно-изоморфне, ако сут изоморфне в обычном смыслу и ако, кроме того, из одговарјаньа $a$ и $\bar a$ следује такоже одговарјанье $\Theta(a)$ и $\Theta(\bar a)$, $H(a)$ и $H(\bar a)$, итд. Ако потом како подгрупы допускајут се само таке, кторе саме допускајут област операторов, тако же група называје се проста, ако не имаје допустимого властного нормалного делительа кроме јединицы, тогда Јордан-Холдеров теорем о композицијном реду оставаје важны в форме: ако група с областју операторов вобче имаје композицијны ред, тогда систем фактор-груп јест в смыслу операторного изоморфизма једнозначно определьен до порјадка. Под тој самоју предпоставкоју важи такоже теорем о разкладу, једнозначном в смыслу операторного изоморфизма, на директно неразложиме факторы.

Вез с проблемом представјень дана јест тым, же к групам с операторами принадлежет особито идеалы и модулы, сматрјане како абелове групы взходно к сложеньу, покы оператори сут даны множеньем с елементами колца. Всако представјенье пораждаје се модулом представјеньа, т. ј. модулом взходно к колцу и телу представјеньа. Точније, всакы клас модулов, т. ј. клас меджу собоју операторно-изоморфных модулов представјеньа, пораждаје то само представјенье. Тако пытанье о једнозначности редукције представјеньа разрешаје се теоремом о композицијном реду и теоремом о директном продукту. Композицијны ред, применьены к модулу представјеньа, даје за всаку матрицу редукцију формы
\[
\begin{pmatrix}
B_1&0&0\\
&B_2&0\\
A_{ik}&\cdots&B_r
\end{pmatrix},
\]
при чем диагоналне матрице, соответны фактор-групам, пораждајут “иредуцибилно” представјенье, т. ј. тако, кторо само уже не допускаје таку редукцију.

Но понеже класы тих фактор-груп сут једнозначно определьене, то само важи за диагоналне представјеньа, једнозначно в смыслу еквиваленције представјень. Аналогично теорем о директном продукту даје једнозначност “разпада” на “неразложиме” составне чести:
\[
\begin{pmatrix}
C_1&&&\\
&C_2&&\\
&&\cdots&\\
&&&C_s
\end{pmatrix},
\]
при чем всако $C_i$ може потом јешче быти једнозначно редуцирано по теорему о композицијном реду.

\emph{Други проблем} хочу начртати в случају представјеньа хиперкомплексных системов, обче -- колц, кторе задовольајут двојно условје цепов. Резултат јест, же всакы иредуцибилны (просты) клас модулов јест клас идеалов; зато композицијне реды из једностранных идеалов чрез своје фактор-групы уже дају все иредуцибилне представјеньа. Точније, достаточне сут композицијне реды в колцу класов остатков по радикалу, кторо ставаје се “пополно редуцибилным” колцом. Проблем јест тым сведены к структурному изследованьу пополно редуцибилных колц; и зато показује се соответным проблему брати представјеньа в телахи автоморфизмов како изходишче. То треба разумети тако: вси прости праве идеалы двустранно просте компоненте такого колца принадлежет к тому самому класу, зато имајут то само колцо автоморфизмов, кторо због простости идеалов ставаје се телом, једновременно колцом автоморфизмов простых левых идеалов. Само просто колцо ставаје се изоморфно систему всих матриц над телом автоморфизмов; и из того представјеньа изходе все представјеньа взходно к подтелу, особито взходно к области коефициентов хиперкомплексного система, когда саме елементы тела автоморфизмов заменьајут се придружеными матрицами. Тако се тут Веддербурнов структурны теорем, ведомы в случају хиперкомплексных системов, ново толкује чрез појем тела автоморфизмов, кторы изходи из обчеј теорије груп, и једновременно познаје се како извор теорије представјень хиперкомплексных системов. Наставајут нове пытаньа Галоисовеј теорије некомутативных тел, кторе тесно связане сут с пытаньем о “разпадных польах”. Једновременно показује се выхлед к структурным теоремам обчих, не пополно редуцибилных колц посредством колц автоморфизмов неразложимых идеалов. Теорија груп, в својем разширјеньу на најобчејше групы с операторами, и тут знову показује се како уредивајучи принцип.

\end{document}

